% English original, revised around the owner's pre-crisis preparation question.
\documentclass[runningheads]{llncs}
\usepackage[T1]{fontenc}
\usepackage{newtxtext}
\usepackage[varvw]{newtxmath}
\usepackage{graphicx,booktabs,amsmath,array,longtable}
\usepackage[hidelinks]{hyperref}
\hypersetup{pdftitle={Supplementary Experimental Record},pdfauthor={},pdfsubject={Anonymous research manuscript}}
\newcommand{\WIP}{\mathrm{WIP}}
\begin{document}
\title{Supplementary Experimental Record}
\titlerunning{Supplementary Evidence}
\author{Anonymous authors}
\authorrunning{Anonymous authors}
\institute{Anonymous submission}
\maketitle
The first sections retain the serial reference experiment P0, its supporting studies and subsequent diagnostics. The final section documents the independent structural extension P1--P6, its complete comparisons and sensitivity analyses. Diagnostics are distinguished from the two frozen confirmatory protocols.

\section{Method: From Operational History to a Prepared Response}
Figure~\ref{fig:workflow} tracks the data used to fit the model, choose a response and test it. Test outcomes evaluate the saved response without changing its selection.
\begin{figure}[t]
\centering\includegraphics[width=\textwidth]{generated/workflow_en.pdf}
\caption{Information flow in the preparation method and its controlled evaluation. Engineers specify process constraints and the assumed effects of actions. The final stage evaluates a policy fixed in the preceding stage.}\label{fig:workflow}
\end{figure}
\subsection{Represent execution and reconstruct observable state}
A case is a request that passes through several operations.

The model $\mathcal M=(N,\mathcal R,K,\alpha,F_D,F_R)$ specifies their order through $N$, groups of interchangeable resources $\mathcal R$, the number available in each group $K$, the assignment of operations to groups $\alpha$, and functions for processing duration and return probability. External arrival schedules are supplied separately. The simulator tracks each case's execution state, occupied resources, queues and the number of unfinished cases $w(t)$, called work in progress, WIP. This count includes cases both waiting and in service. With parallel operations, counting each case once differs from summing activity queues and occupancies.

Ready operations wait for a resource in order of arrival. In the Petri-net representation, tokens record which operations are eligible; input tokens are reserved when an operation joins a queue, and output tokens appear on completion. Occupancy satisfies $0\leq b_r(t)\leq K_r$. Processing duration is sampled once at service start $t_s$, using only the state observable then:
\begin{equation}\label{eq:duration}
\log D_a=f_a(x,\ell(t_s))+\epsilon_a,\qquad \epsilon_a\sim\mathcal N(0,\sigma^2).
\end{equation}
At the modeled return decision at completion $t_c$,
\begin{equation}\label{eq:return}
R\sim\mathrm{Bernoulli}\!\left(\operatorname{sigmoid}(g(x,\ell(t_c)))\right).
\end{equation}
Case attributes $x$ are available at the decision time. The queue length at activity $a$ is $q_a$; workload features include activity-specific $\log(1+q_a/K_{\alpha(a)})$ and $\log(1+w/6)$. The start queue includes the operation about to start; the return queue excludes a prospective returning operation. Future completion states are never start-time features.

We fit ridge regression to log duration ($\alpha=1$) and logistic regression to returns ($C=10$). Static candidates retain case context and finite capacity but omit workload features. Load-dependent candidates retain both workload terms. In the serial implementation, the return feature observes B's queue at C completion and uses its recorded fixed scaling factor 2; duration at B uses capacity 3. These denominators retain their training values when an intervention changes the actual number of workers. The saved implementation fixes both the residual standard deviation and the feature definitions.

An optional selector uses validation-loss gains $\Delta_D$ and $\Delta_R$: it retains duration load terms when $\Delta_D>0.01$ and return load terms when $\Delta_R>0.005$. Duration loss is negative log likelihood on log duration; return loss is binary cross entropy. This selector tests predictive value in the observed validation range. The prospective experiment compares it with unconditional load modeling instead of assuming that predictive selection is optimal for response preparation.

\subsection{Explore stresses and choose a response}
An engineer specifies possible load scenarios $\mathcal S_p$, available actions $\mathcal A$ and their costs. Scenarios can vary arrival intensity, staffing or case mix; the primary experiment tests changes in arrival intensity. Each candidate response pairs an action $a$ with a threshold $q$ of the observed unfinished-case count. The intervention starts after two successive threshold crossings at five-unit samples:
\begin{equation}\label{eq:trigger}
\tau_q=\inf\{t_k:w(t_{k-1})\geq q\ \land\ w(t_k)\geq q\}.
\end{equation}
No realized future arrival schedule is available to this rule. Two samples suppress isolated crossings; this is an explicit queue rule, not a learned crisis classifier. Preparation precedes the disturbance; activation may follow the beginning of a surge while still preceding severe congestion. We do not predict the onset of an unobserved arrival shock.

We evaluate no response, adding one B worker, and halving future return probability. A capacity increase starts additional eligible service without changing durations already sampled. Reduced returns affect only subsequent return decisions. For horizon $H$,
\begin{equation}\label{eq:cost}
J(\pi)=\int_0^H w_\pi(t)\,dt+c_a(H-\tau_q)_+,
\end{equation}
with zero action cost when no trigger occurs. The backlog integral measures unfinished-case time; the implementation approximates it by the trapezoidal rule on the five-unit samples. Costs $c_{\mathrm{capacity}}=2$ and $c_{\mathrm{rework}}=0.5$ are abstract case-time equivalents, not monetary estimates. Every planner receives the same catalog, thresholds, and costs.

Each fitted model simulates the planning library and chooses
\begin{equation}\label{eq:policy}
\widehat\pi=\arg\min_{\pi\in\Pi}\frac{1}{|\mathcal S_p|L}
\sum_{s\in\mathcal S_p}\sum_{l=1}^{L}J_{\widehat{\mathcal M},s,l}(\pi).
\end{equation}
Ties follow a fixed catalog order. Shared random streams reduce incidental differences across models and actions. The selected policy is saved before any test trajectory is generated. Deployment of the rule and evaluation of its effect are subsequent stages.

\subsection{Check coverage and validate the decision}
Before relying on a prepared response, the engineer checks the observed load range, validation gains, resource assumptions and the assumed effects of actions. If two behavioral mechanisms agree on every observed state but differ outside that support, the history alone cannot distinguish their unseen responses. Extrapolation therefore requires assumptions or additional measurements. Modeling load explicitly does not remove this identification limit.

We retain every calibration trajectory and check it for sustained severe congestion. In the prospective experiment a severe event means $w(t)\geq100$ for 200 consecutive time units at the sampling resolution. Trigger lead time is measured against the matched no-action event onset. A no-shock activation is reported separately from a false alarm: baseline operation can itself be unstable. A useful warning need not prevent an event, so realized cost and event incidence are also measured.

The output consists of the saved model, the load scenarios considered, a prepared action and activation threshold, and evidence about their tested operating range. Application in a working organization would also require verifying the effects of interventions and staffing constraints, followed by a controlled deployment assessment. The assumed intervention effects in our simulations do not provide that evidence.

\section{Evaluation Design}
The research method combines artifact construction, controlled simulations, structural discovery, and chronological prediction on public logs. The prospective experiment is primary; earlier campaigns supply supporting evidence. Independent training datasets, rather than repeated trajectories of one fitted model, are the units of uncertainty.

\subsection{Controlled process and prospective experiment}
The process is A$\rightarrow$B$\rightarrow$C, then exit or return to B. Capacities are $(2,3,2)$; baseline mean durations are $(1,2.5,1.4)$. Poisson arrivals carry complexity $c\sim\mathrm{Bernoulli}(0.35)$, multiplying mean service by $1+0.35c$. Lognormal noise has log standard deviation 0.25. Four worlds activate neither feedback (fixed), service feedback, routing feedback, or both. Active mechanisms are
\begin{align}\label{eq:truth}
\mathbb E[D_B\mid q,c]&=2.5(1+0.35c)\left(1+\frac{1.2q}{q+6}\right),\\
P(R=1\mid w,c)&=0.04+0.03c+\frac{0.45z^2}{1+z^2},\quad z=\frac{\max(0,w-6)}{18}.
\end{align}
Inactive feedback removes its load term. Learned log-linear and logistic functions differ from these generating forms. Model simulations clip log duration to $[-4,5]$ and return probability to $[0.001,0.95]$; the supplied truth is not replaced by fitted predictions.

The new experiment, E7, uses 12 independent outer seeds in each world. Eight calibration episodes use rates $\{0.25,0.40,0.55,0.65\}$ twice; three validation episodes use $\{0.30,0.50,0.60\}$. Arrivals end at time 300 and every episode drains fully. All episodes remain in the evaluation; a claim of pre-crisis calibration requires the severe-event audit to pass in the reported stratum.

The planning library combines baseline rates 0.40/0.50 with no pulse, rate 1.10 during $[200,300)$, or rate 1.50 during $[200,400)$. Horizon is 1,400 with two repetitions. Candidate thresholds are 20, 40, and 80. Each model chooses one policy across the six conditions. Test scenarios use new baselines 0.45/0.55 and no pulse, rate 1.35 during $[350,500)$, or rate 1.65 during $[350,600)$, with horizon 1,800 and eight repetitions. Thus, pulse timing and magnitude differ from preparation. Calibration, planning, and test random streams are disjoint. Counter-based operation randomness pairs actions without depending on their event ordering.

The primary candidate is unconditional load modeling. We retain static planning and the validation-selected model, plus no action, capacity-at-20, and rework-at-20 heuristics. The latter two give a strong check on whether learning which action to take adds value. A development pilot used one separate seed per world and is excluded from inference. Its inspection preceded protocol freezing; it motivated the unconditional-load primary candidate and common planning noise. Production includes 48 fitted datasets, 12,096 planning rollouts and 16,128 test rollouts. No production outcome changes model selection or the protocol.

The primary endpoint is realized cost averaged equally over six test conditions and their eight repetitions. Routing is the primary mechanism, specified from the preceding experiments. The two primary contrasts compare load planning with static planning and capacity-at-20, using paired outer-seed differences and two-sided 97.5\% Student-$t$ intervals with 11 degrees of freedom, a Bonferroni adjustment for two claims. Rework-at-20 is always reported, including ties. Other 95\% intervals are descriptive. A secondary cross-world comparison averages each world's cost normalized by its no-action cost within each seed, giving mechanisms equal weight. Sensitivity checks use cost multipliers 0.5, 1, and 2; policies are reselected using saved planning results only, then scored on saved test actions.

\subsection{Supporting simulation and real-log checks}
Earlier fixed-structure fits use six seeds per world, training arrival rates 0.40, 0.65, 0.80, 1.05, and 1.30 and validation rates 0.50, 0.90, and 1.15, with arrivals to time 400 and complete drain. An initial audit (E1) checks 235,114 operations and 432 simulations. E5 reuses those fits with independent endogenous noise, baseline rates 0.40/0.50/0.65, optional rate-1.50 pulses during $[300,500)$, horizon 2,000, and eight repetitions. Its stronger empirical static baseline resamples log durations within each activity and complexity class. A two-layer, 32-unit tanh neural comparator uses the same context/load inputs. E5 also tests B+1 and halved returns at baseline 0.50, both applied from time zero; these are not E7's triggered responses. E5 contains 7,680 new test simulations.

The structural campaign (E2/E4) has serial and parallel topologies, four worlds, and six seeds. The parallel truth has A, concurrent B/C with a shared three-worker pool, an AND-join, D, and repetition of the joint block. Complexity is absent. Training rates are 0.25/0.40/0.60/0.80 and validation rates 0.35/0.70, arrivals to time 300 followed by drain. Inductive Miner in PM4Py 2.7.23.12 discovers structure; overlapping observed worker identities identify shared pools and distinct workers their capacities. A compositional timed executor supports sequence, parallel and binary-loop blocks. One known return function is reused at discovered loop locations, a restriction relevant to overgeneralization. Prosimos 2.1.0 executes discovered BPMN with matched static parameters. Serial/parallel test baselines are 0.50/0.40, pulse rate 1.20 during $[200,400)$, horizon 1,200, six repetitions: 3,456 new simulations.

For these earlier trajectory tests, error is the root mean square difference between candidate and truth mean WIP paths, divided by $\max(1,\text{mean truth WIP})$. It is calculated per outer seed, then averaged. Paired differences use 10,000-resample percentile bootstrap intervals across six seeds, descriptive and unadjusted. These scores assess trajectory reconstruction, not realized intervention benefit.

The real-log study uses original version-1 BPI 2012 (13,087 cases; 262,200 events) and BPI 2017 (31,509 cases; 1,202,267 events)~\cite{bpi12,bpi17}. Lifecycle start/resume intervals are paired with observed stopping events; they are wall-clock intervals, not direct labor measurements. Chronological case cohorts split 60/20/20, with labels crossing boundaries purged. Activity, start-time resource, case context, elapsed history, and calendar form the context baseline. Load adds active intervals, open work items and recent work starts globally and by activity/resource. Unknown start resources remain unknown; replay is prefix-invariant. Histogram gradient boosting is the strong reported comparator. Targets are $\log(1+\mathrm{seconds})$ duration and next observed work start conditional on continuation, not an exit decision. Paired loss differences use 2,000 weekly-block bootstrap resamples, six and sixteen test weeks. This is predictive association, not real-world policy evaluation.

\section{Results}
\subsection{Preparing and activating responses before a crisis}
All 48 production tasks pass coverage, completion-hash and raw-path audits. None of the 528 calibration/validation episodes meets the severe-event criterion. Maximum calibration WIP is 16, 61, 19 and 111 for fixed, service, routing and combined feedback, respectively. The last value includes a brief excursion above 100, so ``pre-crisis'' means absence of the defined sustained event, not uniformly low load. No episode is removed.

Table~\ref{tab:e7cost} reports all mechanisms and comparators. In the primary routing setting, load planning chooses rework-at-20 in 9 of 12 fits and capacity-at-20 in the other 3. Mean realized cost falls from 92,832.6 with static planning to 57,351.8, a 38.2\% reduction. The capacity-at-20 comparison gives a 38.1\% reduction. Both adjusted intervals in Table~\ref{tab:e7contrast} exclude zero, satisfying the frozen primary criterion. Rework-at-20 alone costs 45,313.5: load planning is numerically worse, and its paired interval includes zero. This is no evidence of equivalence or superiority over that heuristic. The predictive selector retains return load terms in only 2 of 12 fits and chooses rework in only 1; its mean cost is 88,834.8.

\begin{table}[t]
\caption{E7 mean realized cost in thousands of case-time-equivalent units, over 12 independent fits per mechanism and all test conditions. ``Selected'' uses validation-based component selection.}\label{tab:e7cost}
\centering\small\setlength{\tabcolsep}{3pt}
\begin{tabular}{lrrrrrr}\toprule
Mechanism & \shortstack{No\\action} & \shortstack{Static\\planning} & \shortstack{Load\\planning} & \shortstack{Selected\\planning} & \shortstack{Capacity\\at 20} & \shortstack{Rework\\at 20}\\\midrule
Fixed & 23.6 & 18.0 & 19.9 & 18.0 & 17.7 & 22.6 \\
Service & 211.7 & 111.0 & 111.0 & 111.0 & 111.0 & 202.7 \\
Routing & 147.8 & 92.8 & 57.4 & 88.8 & 92.6 & 45.3 \\
Combined & 375.2 & 309.2 & 288.3 & 292.5 & 309.2 & 284.1 \\
\bottomrule

\end{tabular}
\end{table}
\begin{table}[t]
\caption{Routing: load planning minus comparator cost, in thousands. Primary intervals are 97.5\%; the rework comparison is descriptive at 95\%. Negative differences favor load planning.}\label{tab:e7contrast}
\centering\small
\begin{tabular}{lrrr}\toprule
Comparator & Difference & Paired interval & Level\\\midrule
Static planning & $-35.48$ & $[-51.56, -19.40]$ & 97.5\% \\
Capacity at 20 & $-35.27$ & $[-51.23, -19.32]$ & 97.5\% \\
Rework at 20 & $12.04$ & $[-1.87, 25.94]$ & 95\% \\
\bottomrule

\end{tabular}
\end{table}

The response is activated before the matched no-action crisis in all 289 at-risk routing trajectories, with mean lead time 108.0 units. It prevents the defined event in 93 of these trajectories; 196 still experience one. Neither an activation nor a crisis occurs in the 192 no-shock routing trajectories. Figure~\ref{fig:e7} illustrates faster recovery and the variation caused by learned policy choices. These are conditional, descriptive warning measurements, not a calibrated crisis-classification score.

\begin{figure}[t]
\centering\includegraphics[width=\textwidth]{generated/e7_response_en.pdf}
\caption{Routing responses under two held-out pulse scenarios. Lines average eight trajectories within each of 12 fits, then average fits. The load-planning band is one standard deviation across fit means, not a confidence interval. Gray regions mark pulses; the horizontal line marks WIP 100. The rework heuristic shows the limit of learned action selection.}\label{fig:e7}
\end{figure}

Benefits vary by mechanism. For service feedback, every planner chooses capacity-at-20 and all costs coincide. With combined feedback, load planning lowers cost by 6.7\% against static/capacity planning, but severe-event incidence rises from 69.3\% to 71.0\%: minimizing backlog cost is not identical to preventing the defined event. Without feedback, load planning raises cost by 10.9\% against static planning, with descriptive paired difference 1,959.3 and interval $[219.5,3699.1]$. Across mechanisms, no-shock activation rates are 0.0\%, 14.1\%, 0.0\% and 28.6\%; matching no-action crisis rates are 0.0\%, 2.1\%, 0.0\% and 19.3\%. Baseline instability therefore explains some, but not all, no-shock activations.

The equal-weight cross-mechanism analysis favors load planning at default costs: normalized differences are $-0.0536$ against static planning, $-0.0503$ against capacity-at-20, and $-0.1137$ against rework-at-20. These are differences of within-world cost ratios, not percentage reductions in pooled raw cost. Table~\ref{tab:e7sensitivity} gives descriptive intervals and all prespecified cost multipliers. At twice the action cost, the static and capacity intervals include zero; the advantage over rework remains marginal. Thus, adaptation across mechanisms is supported for the stated decision problem, not independently of its cost assumptions.

\begin{table}[t]
\caption{Cross-mechanism normalized cost differences, load planning minus comparator, with descriptive 95\% intervals. Policies are reselected from planning costs at each multiplier.}\label{tab:e7sensitivity}
\centering\small\setlength{\tabcolsep}{3pt}
\begin{tabular}{lrrr}\toprule
Cost multiplier & Static planning & Capacity at 20 & Rework at 20\\\midrule
0.5 & \shortstack{$-0.0467$\\$[-0.0772,-0.0161]$} & \shortstack{$-0.0467$\\$[-0.0772,-0.0161]$} & \shortstack{$-0.1201$\\$[-0.1501,-0.0901]$} \\
1.0 & \shortstack{$-0.0536$\\$[-0.0824,-0.0247]$} & \shortstack{$-0.0503$\\$[-0.0790,-0.0217]$} & \shortstack{$-0.1137$\\$[-0.1420,-0.0855]$} \\
2.0 & \shortstack{$-0.0136$\\$[-0.0618,0.0346]$} & \shortstack{$-0.0037$\\$[-0.0510,0.0436]$} & \shortstack{$-0.0471$\\$[-0.0939,-0.0002]$} \\
\bottomrule

\end{tabular}
\end{table}


\subsection{Supporting evidence and transfer failures}
The earlier broad-load E1 fits select the correct pair of feedback components in all 24 datasets. E5's independent-noise comparison improves on an empirical static baseline in three mechanisms at baseline 0.65 with a pulse (Table~\ref{tab:e5}). Mean-path error reductions are 76.2\%, 88.1\%, and 64.4\%; all three paired descriptive intervals lie below zero. The fixed-world difference includes zero. These numbers concern trajectory error in the stated setting, not E7 cost or a universal ranking.
\begin{table}[t]
\caption{E5 pulse at baseline 0.65. Mean WIP-path NRMSE over six outer seeds. Difference: selected minus empirical static, with descriptive 95\% interval.}\label{tab:e5}
\centering\small\setlength{\tabcolsep}{4pt}
\begin{tabular}{lrrrrl}\toprule
Mechanism & Empirical & Selected & Neural & Difference & Interval\\\midrule
Fixed & 0.0942 & 0.1067 & 0.2044 & $+0.0126$ & $[-0.0084,+0.0381]$ \\
Service & 0.1745 & 0.0416 & 0.0499 & $-0.1329$ & $[-0.1684,-0.1000]$ \\
Routing & 0.7733 & 0.0922 & 0.0750 & $-0.6811$ & $[-0.7488,-0.6118]$ \\
Combined & 0.0920 & 0.0328 & 0.0475 & $-0.0592$ & $[-0.0693,-0.0496]$ \\
\bottomrule

\end{tabular}
\end{table}

At baseline 0.50, routing truth ends with 1.46 unfinished cases under constant arrivals and 109.77 after the pulse, although arrivals returned to baseline 1,500 time units earlier. Service feedback also yields persistent backlog. Combined feedback can be unstable even without a pulse. The observation supports a finite-horizon traffic-jam effect; it is not a stationary phase-transition result.

With discovered structure, selected behavior improves on matched Prosimos for serial service, serial combined and parallel service feedback (Table~\ref{tab:structure}), with paired NRMSE differences $-0.1987$, $-0.7166$ and $-0.0998$ and descriptive intervals excluding zero. Parallel routing and combined feedback instead worsen by 0.2189 and 0.2019. Discovery allows independent repetitions where the truth repeats a joint block. Token fitness is 1.0 for both topologies, but parallel precision averages 0.4945. Thus, high fitness does not ensure that discovered loops preserve system load. Serial routing also selects static return behavior in every seed under its narrower observation range.
\begin{table}[t]
\caption{All discovered-structure pulse settings, mean NRMSE over six seeds. D: discovered structure; G: supplied structure.}\label{tab:structure}
\centering\small\setlength{\tabcolsep}{4pt}
\begin{tabular}{lrrrr}\toprule
Topology / mechanism & D static & D selected & Prosimos & G selected\\\midrule
Serial / Fixed & 0.1983 & 0.1874 & 0.2133 & 0.1850 \\
Serial / Service & 0.9207 & 0.7219 & 0.9206 & 0.7119 \\
Serial / Routing & 1.2318 & 1.2295 & 1.2351 & 1.2318 \\
Serial / Combined & 0.9779 & 0.2635 & 0.9801 & 0.2743 \\
Parallel / Fixed & 0.0768 & 0.0914 & 0.0609 & 0.0800 \\
Parallel / Service & 0.1968 & 0.0904 & 0.1901 & 0.0703 \\
Parallel / Routing & 0.6397 & 0.8500 & 0.6312 & 0.4988 \\
Parallel / Combined & 0.1071 & 0.3141 & 0.1123 & 0.1559 \\
\bottomrule

\end{tabular}
\end{table}

On BPI 2017, boosted duration MAE decreases from 1.188915 to 1.179020 with load, a relative improvement of 0.83\% and paired interval $[-0.014642,-0.006249]$. Table~\ref{tab:real} reports all log/target pairs. BPI 2012 duration is essentially unchanged, its next-start loss worsens, and the BPI 2017 next-start interval includes zero. Load information is useful for a particular target; these observational checks do not validate a real intervention.
\begin{table}[t]
\caption{All boosted real-log tests. Duration: MAE on log(1+seconds); next start: log loss. Intervals describe load-minus-context differences.}\label{tab:real}
\centering\small\setlength{\tabcolsep}{4pt}
\begin{tabular}{lrrrr}\toprule
Log / target & Test rows & Context & With load & Difference interval\\\midrule
BPI12 / Duration & 12,653 & 1.0308 & 1.0308 & $[-0.0051,+0.0033]$ \\
BPI12 / Next start & 10,781 & 0.3585 & 0.3792 & $[+0.0022,+0.0436]$ \\
BPI17 / Duration & 47,078 & 1.1889 & 1.1790 & $[-0.0146,-0.0062]$ \\
BPI17 / Next start & 26,101 & 0.0815 & 0.0838 & $[-0.0005,+0.0058]$ \\
\bottomrule

\end{tabular}
\end{table}

Finally, earlier observational accuracy fails an intervention-ranking check. In the combined world, true final WIP with pulse/B+1/halved returns is 628.50/473.52/414.25, but the selected model predicts 654.17/230.67/425.38 and prefers the wrong action. A post-result coordinate diagnostic (E6) keeps B's duration-feature denominator at training capacity while increasing physical capacity. Its prediction moves to 413.81 and error falls from about 0.568 to 0.180, but the ranking remains wrong. This diagnostic adds 384 simulations and is not an independently confirmed repair. E7 adopts fixed feature coordinates prospectively and tests the actual selected response instead of treating observational fit as sufficient.

\section{Discussion and Limitations}
The intended IS engineering contribution is a preparation procedure with explicit evidence at each stage: operational-state reconstruction, behavioral modeling, hypothetical stress testing, response selection, and independent evaluation. It connects discovery to design and operation of process-aware systems. A response card can specify its observable trigger, resource requirements, assumed action effect, estimated cost and tested load range. Our artifact supplies the computational parts of such a card; organizational feasibility and deployment remain outside the present evaluation.

The scope of inference matters. E7 supplies serial structure and intervention semantics. The earlier supporting experiments test structural discovery without reproducing the complete response pipeline. The separate structural extension reported below evaluates that pipeline on six more complex structures. Real-log tests concern association under imperfect lifecycle and staffing information. Neither synthetic success nor the modest real-log predictive gain identifies a causal intervention effect in an organization. Human approval, resource availability, intervention delay and compliance constraints would need explicit modeling in deployment.

The load range is another boundary. Calm validation can support short-term prediction while failing to reveal a dependency relevant to a future shock. Conversely, an unconstrained load model can extrapolate noise into an implausible crisis. The calibration audit and the static/selected comparisons expose this trade-off; they do not solve general out-of-distribution identification. Additional controlled measurements or safe stress probes may be necessary when action rankings depend on unobserved states.

Twelve outer seeds still give coarse uncertainty, and the four synthetic worlds are not a population sample of organizations. Primary intervals adjust two prespecified comparisons; remaining intervals are descriptive. The equal-weight scenario library and abstract action costs encode a decision problem, not observed incident frequencies or prices. We report cost sensitivity and simple-policy ties to avoid overstating the value of learning. Five-unit observations, a fixed 200-unit crisis duration, known action efficacy and a finite test horizon limit interpretations of prevention and lead time. Permanent stability, rare-event reliability and generalization to missing events or changing case mixes are not established.

The earlier campaigns contain 11,952 newly generated test simulations, separate from E7 and its excluded pilot. Simulation counts describe computational coverage; uncertainty uses independent fits or chronological blocks. All adverse structural and real-log comparisons remain visible. We do not claim superiority over unexecuted full simulation-discovery systems or causal prescriptive-monitoring methods.

\section{Additional Diagnostic Record}
The following analyses were specified after the primary results were known. Original policies remain frozen; all 48 fits and all declared diagnostic cells are retained. Intervals are descriptive 95\% paired intervals over independent training fits.

\begin{table}[p]\caption{Mean cost in thousands for every mechanism and diagnostic comparator.}\centering\small\setlength{\tabcolsep}{3pt}\begin{tabular}{p{.40\textwidth}rrrr}\toprule
Method & Fixed & Service & Routing & Combined\\\midrule
No action & 23.6 & 211.7 & 147.8 & 375.2\\
Static planning & 18.0 & 111.0 & 92.8 & 309.2\\
Predictive selection & 18.0 & 111.0 & 88.8 & 292.5\\
Feedback-strength rule & 18.9 & 111.0 & 73.0 & 309.2\\
Retained-load planning & 19.9 & 111.0 & 57.4 & 288.3\\
Capacity at 20 & 17.7 & 111.0 & 92.6 & 309.2\\
Rework at 20 & 22.6 & 202.7 & 45.3 & 284.1\\
Known-generator planner & 17.8 & 111.0 & 45.3 & 284.1\\
Hindsight envelope & 17.7 & 111.0 & 45.3 & 284.1\\
\bottomrule\end{tabular}\end{table}
\begin{table}[p]\caption{Cost above the known-generator planner: mean and 95\% interval, in thousands. Negative values are possible because the oracle uses finite planning samples and different planning scenarios.}\centering\small\setlength{\tabcolsep}{3pt}\begin{tabular}{llr}\toprule
Mechanism & Method & Excess cost [interval]\\\midrule
Fixed & Static planning & 0.17 [-0.08, 0.42]\\
Fixed & Predictive selection & 0.17 [-0.08, 0.42]\\
Fixed & Feedback-strength rule & 1.09 [-0.19, 2.36]\\
Fixed & Retained-load planning & 2.13 [0.41, 3.85]\\
Service & Static planning & 0.00 [0.00, 0.00]\\
Service & Predictive selection & 0.00 [0.00, 0.00]\\
Service & Feedback-strength rule & 0.00 [0.00, 0.00]\\
Service & Retained-load planning & 0.00 [0.00, 0.00]\\
Routing & Static planning & 47.52 [45.70, 49.34]\\
Routing & Predictive selection & 43.52 [34.63, 52.42]\\
Routing & Feedback-strength rule & 27.66 [12.12, 43.20]\\
Routing & Retained-load planning & 12.04 [-1.87, 25.94]\\
Combined & Static planning & 25.05 [23.43, 26.68]\\
Combined & Predictive selection & 8.32 [0.44, 16.19]\\
Combined & Feedback-strength rule & 25.05 [23.43, 26.68]\\
Combined & Retained-load planning & 4.20 [-2.05, 10.45]\\
\bottomrule\end{tabular}\end{table}
\begin{table}[p]\caption{Intervention efficacy: frozen decisions, assumed reduction 50\%; actual reductions shown. Mean costs in thousands.}\centering\small\setlength{\tabcolsep}{3pt}\begin{tabular}{llrrrrr}\toprule
Mechanism & Effect & Static & Selected & Load & Capacity & Rework\\\midrule
Fixed & 30\% & 18.0 & 18.0 & 20.1 & 17.7 & 23.2\\
Fixed & 50\% & 18.0 & 18.0 & 19.9 & 17.7 & 22.6\\
Fixed & 70\% & 18.0 & 18.0 & 19.7 & 17.7 & 22.0\\
Service & 30\% & 111.0 & 111.0 & 111.0 & 111.0 & 206.5\\
Service & 50\% & 111.0 & 111.0 & 111.0 & 111.0 & 202.7\\
Service & 70\% & 111.0 & 111.0 & 111.0 & 111.0 & 199.1\\
Routing & 30\% & 92.8 & 91.1 & 77.1 & 92.6 & 71.7\\
Routing & 50\% & 92.8 & 88.8 & 57.4 & 92.6 & 45.3\\
Routing & 70\% & 92.8 & 87.8 & 47.6 & 92.6 & 32.2\\
Combined & 30\% & 309.2 & 316.8 & 319.2 & 309.2 & 321.1\\
Combined & 50\% & 309.2 & 292.5 & 288.3 & 309.2 & 284.1\\
Combined & 70\% & 309.2 & 268.1 & 258.0 & 309.2 & 248.1\\
\bottomrule\end{tabular}\end{table}
\begin{table}[p]\caption{Mixture sensitivity: normalized load-minus-comparator mean difference; negative favors retained-load planning. The fixed-world weight varies and the other three worlds share the remaining weight equally.}\centering\small\setlength{\tabcolsep}{3pt}\begin{tabular}{p{.37\textwidth}rrrr}\toprule
Comparator & 0 & 0.25 & 0.5 & 1\\\midrule
Static planning & -0.099 & -0.054 & -0.008 & 0.083\\
Predictive selection & -0.075 & -0.036 & 0.004 & 0.083\\
Capacity at 20 & -0.099 & -0.050 & -0.002 & 0.095\\
Rework at 20 & -0.114 & -0.114 & -0.113 & -0.113\\
Feedback-strength rule & -0.054 & -0.030 & -0.005 & 0.044\\
Known-generator planner & 0.030 & 0.045 & 0.060 & 0.090\\
\bottomrule\end{tabular}\end{table}
\begin{table}[p]\caption{Routing event sensitivity: percent of trajectories reaching the sustained-event criterion; original frozen policies.}\centering\small\setlength{\tabcolsep}{3pt}\begin{tabular}{rrrrrr}\toprule
WIP & Duration & No action & Static & Load & Rework\\\midrule
80 & 100 & 61.6 & 50.5 & 39.9 & 37.5\\
80 & 200 & 59.0 & 43.6 & 37.3 & 34.2\\
80 & 300 & 55.6 & 38.2 & 34.5 & 33.3\\
100 & 100 & 54.2 & 39.2 & 35.6 & 33.7\\
100 & 200 & 50.2 & 35.8 & 34.0 & 33.3\\
100 & 300 & 46.9 & 34.4 & 33.3 & 32.5\\
120 & 100 & 43.9 & 34.4 & 34.0 & 33.3\\
120 & 200 & 42.7 & 34.0 & 33.2 & 32.1\\
120 & 300 & 41.3 & 33.9 & 30.4 & 29.0\\
\bottomrule\end{tabular}\end{table}
\begin{table}[p]\caption{Routing cost sensitivity: horizon and terminal unfinished-case weight. Mean costs in thousands; no response is selected again.}\centering\small\setlength{\tabcolsep}{3pt}\begin{tabular}{rrrrrr}\toprule
Horizon & Weight & Static & Selected & Load & Rework\\\midrule
1400 & 0 & 77.0 & 74.2 & 52.2 & 43.7\\
1400 & 100 & 81.9 & 78.7 & 53.8 & 44.3\\
1400 & 200 & 86.8 & 83.2 & 55.4 & 44.9\\
1600 & 0 & 85.9 & 82.4 & 55.1 & 44.6\\
1600 & 100 & 89.7 & 85.9 & 56.3 & 44.9\\
1600 & 200 & 93.5 & 89.5 & 57.4 & 45.3\\
1800 & 0 & 92.8 & 88.8 & 57.4 & 45.3\\
1800 & 100 & 95.8 & 91.6 & 58.4 & 45.6\\
1800 & 200 & 98.7 & 94.3 & 59.4 & 45.9\\
\bottomrule\end{tabular}\end{table}
Complete per-fit choices, all threshold counts, raw efficacy and sensitivity cells, discovered nets, compiler proofs, and the machine-readable response card are provided in the archive. 4032 new oracle planning rollouts; 9156 new effect rollouts and 4668 exactly reused non-activation trajectories. Reuse is valid because efficacy changes only events after activation, which never occurs on those paths.
\clearpage
\begin{table}[p]\caption{Activation and prevention counts under frozen load-planning responses. The last column counts activations on trajectories with no severe event in the matched no-action run; this is distinct from calm-control activation.}\centering\small\setlength{\tabcolsep}{3pt}\begin{tabular}{lrrrr}\toprule
Mechanism & At risk & Early & Prevented & \shortstack{Activation without\\baseline event}\\\midrule
Fixed & 117 & 109 & 74 & 239\\
Service & 382 & 382 & 138 & 29\\
Routing & 289 & 289 & 93 & 95\\
Combined & 421 & 421 & 12 & 18\\
\bottomrule\end{tabular}\end{table}\clearpage

\clearpage

\section{Independent Structural Extension}

This extension contains 288 independent training datasets and 3168 audited calibration/validation episodes. All datasets satisfy the declared pre-crisis criterion. Known and discovered structures share learned behavioral parameters and supplied resource maps. No discovery repair or test-based policy reselection is applied. The complete machine-readable comparison, including every fixed action and component ablation, is in the archive file \texttt{analysis/full\_summary.csv}.

\subsection{Costs across all structures and mechanisms}

Tables report the mean of per-fit cost ratios to matched no action, which is 1. Regret is retained-load minus generator-planner cost on that scale. Severe-event incidence refers to retained-load. K denotes known and D discovered structure. Every row has 12 independent fits; the heuristic and generator columns use the same comparator in both structure variants.

\begin{table}[htbp]\caption{P1 normalized costs and event incidence.}\centering\small\setlength{\tabcolsep}{3pt}

\begin{tabular}{llrrrrrrr}\toprule

Model & Mechanism & Static & Pred. & Ret. & Heur. & Gen. & Regret & Event\\\midrule

K & fixed & 0.817 & 0.800 & 0.807 & 0.801 & 0.798 & 0.009 & 0.082\\

K & service & 0.599 & 0.630 & 0.639 & 0.670 & 0.578 & 0.061 & 0.399\\

K & routing & 0.587 & 0.549 & 0.502 & 0.501 & 0.429 & 0.073 & 0.365\\

K & combined & 0.729 & 0.692 & 0.624 & 0.636 & 0.588 & 0.036 & 0.641\\

D & fixed & 0.760 & 0.760 & 0.760 & 0.801 & 0.798 & -0.038 & 0.035\\

D & service & 0.580 & 0.580 & 0.577 & 0.670 & 0.578 & -0.001 & 0.366\\

D & routing & 0.575 & 0.575 & 0.568 & 0.501 & 0.429 & 0.139 & 0.378\\

D & combined & 0.722 & 0.722 & 0.721 & 0.636 & 0.588 & 0.133 & 0.648\\

\bottomrule\end{tabular}\end{table}

\begin{table}[htbp]\caption{P2 normalized costs and event incidence.}\centering\small\setlength{\tabcolsep}{3pt}

\begin{tabular}{llrrrrrrr}\toprule

Model & Mechanism & Static & Pred. & Ret. & Heur. & Gen. & Regret & Event\\\midrule

K & fixed & 0.955 & 0.955 & 0.988 & 1.077 & 0.953 & 0.035 & 0.102\\

K & service & 0.647 & 0.647 & 0.701 & 0.820 & 0.630 & 0.072 & 0.333\\

K & routing & 0.731 & 0.614 & 0.428 & 0.514 & 0.428 & 0.000 & 0.325\\

K & combined & 0.751 & 0.640 & 0.582 & 0.605 & 0.521 & 0.061 & 0.514\\

D & fixed & 0.918 & 0.918 & 1.072 & 1.077 & 0.953 & 0.119 & 0.095\\

D & service & 0.637 & 0.637 & 0.878 & 0.820 & 0.630 & 0.249 & 0.351\\

D & routing & 0.772 & 0.772 & 0.897 & 0.514 & 0.428 & 0.469 & 0.370\\

D & combined & 0.748 & 0.751 & 0.798 & 0.605 & 0.521 & 0.278 & 0.576\\

\bottomrule\end{tabular}\end{table}

\begin{table}[htbp]\caption{P3 normalized costs and event incidence.}\centering\small\setlength{\tabcolsep}{3pt}

\begin{tabular}{llrrrrrrr}\toprule

Model & Mechanism & Static & Pred. & Ret. & Heur. & Gen. & Regret & Event\\\midrule

K & fixed & 0.930 & 0.930 & 0.958 & 0.920 & 0.923 & 0.036 & 0.313\\

K & service & 0.535 & 0.535 & 0.643 & 0.535 & 0.535 & 0.108 & 0.535\\

K & routing & 0.630 & 0.630 & 0.650 & 0.630 & 0.516 & 0.133 & 0.651\\

K & combined & 0.878 & 0.878 & 0.871 & 0.880 & 0.839 & 0.032 & 0.667\\

D & fixed & 0.930 & 0.930 & 0.958 & 0.920 & 0.923 & 0.036 & 0.313\\

D & service & 0.535 & 0.535 & 0.643 & 0.535 & 0.535 & 0.108 & 0.535\\

D & routing & 0.630 & 0.630 & 0.650 & 0.630 & 0.516 & 0.133 & 0.651\\

D & combined & 0.878 & 0.878 & 0.871 & 0.880 & 0.839 & 0.032 & 0.667\\

\bottomrule\end{tabular}\end{table}

\begin{table}[htbp]\caption{P4 normalized costs and event incidence.}\centering\small\setlength{\tabcolsep}{3pt}

\begin{tabular}{llrrrrrrr}\toprule

Model & Mechanism & Static & Pred. & Ret. & Heur. & Gen. & Regret & Event\\\midrule

K & fixed & 0.720 & 0.720 & 0.805 & 0.888 & 0.720 & 0.085 & 0.255\\

K & service & 0.545 & 0.545 & 0.649 & 0.649 & 0.544 & 0.105 & 0.564\\

K & routing & 0.520 & 0.528 & 0.610 & 0.620 & 0.512 & 0.098 & 0.490\\

K & combined & 0.801 & 0.809 & 0.828 & 0.839 & 0.800 & 0.029 & 0.672\\

D & fixed & 0.720 & 0.720 & 0.805 & 0.888 & 0.720 & 0.085 & 0.255\\

D & service & 0.545 & 0.545 & 0.649 & 0.649 & 0.544 & 0.105 & 0.564\\

D & routing & 0.520 & 0.528 & 0.610 & 0.620 & 0.512 & 0.098 & 0.490\\

D & combined & 0.801 & 0.809 & 0.828 & 0.839 & 0.800 & 0.029 & 0.672\\

\bottomrule\end{tabular}\end{table}

\begin{table}[htbp]\caption{P5 normalized costs and event incidence.}\centering\small\setlength{\tabcolsep}{3pt}

\begin{tabular}{llrrrrrrr}\toprule

Model & Mechanism & Static & Pred. & Ret. & Heur. & Gen. & Regret & Event\\\midrule

K & fixed & 0.799 & 0.799 & 0.950 & 0.985 & 0.800 & 0.150 & 0.306\\

K & service & 0.348 & 0.350 & 0.352 & 0.460 & 0.348 & 0.003 & 0.483\\

K & routing & 1.030 & 0.727 & 0.473 & 0.422 & 0.422 & 0.050 & 0.354\\

K & combined & 1.142 & 1.022 & 1.008 & 0.997 & 0.984 & 0.025 & 0.672\\

D & fixed & 0.799 & 0.799 & 0.950 & 0.985 & 0.800 & 0.150 & 0.306\\

D & service & 0.348 & 0.350 & 0.352 & 0.460 & 0.348 & 0.003 & 0.483\\

D & routing & 1.030 & 0.727 & 0.473 & 0.422 & 0.422 & 0.050 & 0.354\\

D & combined & 1.142 & 1.022 & 1.008 & 0.997 & 0.984 & 0.025 & 0.672\\

\bottomrule\end{tabular}\end{table}

\begin{table}[htbp]\caption{P6 normalized costs and event incidence.}\centering\small\setlength{\tabcolsep}{3pt}

\begin{tabular}{llrrrrrrr}\toprule

Model & Mechanism & Static & Pred. & Ret. & Heur. & Gen. & Regret & Event\\\midrule

K & fixed & 0.753 & 0.753 & 0.799 & 1.051 & 0.753 & 0.045 & 0.217\\

K & service & 0.494 & 0.494 & 0.633 & 0.920 & 0.492 & 0.140 & 0.439\\

K & routing & 0.575 & 0.575 & 0.422 & 0.460 & 0.349 & 0.073 & 0.401\\

K & combined & 0.745 & 0.693 & 0.602 & 0.685 & 0.590 & 0.012 & 0.665\\

D & fixed & 0.740 & 0.740 & 0.740 & 1.051 & 0.753 & -0.014 & 0.174\\

D & service & 0.492 & 0.492 & 0.492 & 0.920 & 0.492 & 0.000 & 0.359\\

D & routing & 0.568 & 0.568 & 0.605 & 0.460 & 0.349 & 0.256 & 0.528\\

D & combined & 0.743 & 0.743 & 0.743 & 0.685 & 0.590 & 0.153 & 0.667\\

\bottomrule\end{tabular}\end{table}

\clearpage

\subsection{Primary contrasts and discovery checks}

The six primary routing contrasts use Bonferroni-adjusted intervals, controlling the family error rate at 0.05. Values below are absolute retained-minus-predictive cost differences divided by 1000, not percent changes. Positive differences favor predictive selection.

\begin{center}\small\begin{tabular}{lrrr}\toprule

Process & Difference & Lower & Upper\\\midrule

P1 & -1.081 & -2.317 & 0.155\\

P2 & -4.694 & -10.293 & 0.905\\

P3 & 2.056 & -9.309 & 13.420\\

P4 & 5.092 & 1.931 & 8.252\\

P5 & -14.544 & -31.191 & 2.103\\

P6 & -11.514 & -25.155 & 2.126\\

\bottomrule\end{tabular}\end{center}

\begin{center}\small\begin{tabular}{lrrrrrr}\toprule

Process & Fitness & Precision & Equal & Added & Missing & Same policy\\\midrule

P1 & 1.000 & 0.497 & 0/48 & 48/48 & 0/48 & 25/48\\

P2 & 1.000 & 0.642 & 0/48 & 48/48 & 0/48 & 3/48\\

P3 & 1.000 & 0.996 & 48/48 & 0/48 & 0/48 & 48/48\\

P4 & 1.000 & 1.000 & 48/48 & 0/48 & 0/48 & 48/48\\

P5 & 1.000 & 1.000 & 48/48 & 0/48 & 0/48 & 48/48\\

P6 & 1.000 & 0.487 & 0/48 & 48/48 & 0/48 & 17/48\\

\bottomrule\end{tabular}\end{center}

Language checks compare visible execution languages with a fixed state bound; unresolved checks are retained in \texttt{discovery\_summary.json}. Same policy counts exact retained-load action/threshold agreement between known and discovered structure across all four mechanisms. The 144 models with added paths are recorded as reference-conformance failures in \texttt{structural\_conformance.json}; their unrepaired costs remain in the error analysis. This reference-based evaluation flag is separate from computational failure and is not a deployed rejection rule.

\subsection{Selector sensitivity}

The sweeps reuse saved planning selections and independent test actions. They are diagnostic; no threshold is selected from test costs. The following tables concern routing feedback and known structure. Duration and return counts give the number of retained components out of 12. Agree is action-type agreement with generator planning; regret is absolute cost difference divided by 1000. Other mechanisms and discovered structures are included in \texttt{selector\_summary.csv}.

\begingroup\small\setlength{\tabcolsep}{3pt}\begin{longtable}{llrrrrrr}\caption{Selector sweep: return.}\\\toprule

Process & Threshold & D & R & Agree & Cost ratio & Regret & Event\\\midrule\endfirsthead

\toprule Process & Threshold & D & R & Agree & Cost ratio & Regret & Event\\\midrule\endhead

P1 & 0 & 0 & 6 & 0.417 & 0.527 & 2.450 & 0.378\\

P1 & 0.001 & 0 & 3 & 0.250 & 0.549 & 2.920 & 0.384\\

P1 & 0.0025 & 0 & 3 & 0.250 & 0.549 & 2.920 & 0.384\\

P1 & 0.005 & 0 & 3 & 0.250 & 0.549 & 2.920 & 0.384\\

P1 & 0.01 & 0 & 2 & 0.167 & 0.561 & 3.186 & 0.385\\

P2 & 0 & 0 & 8 & 0.667 & 0.566 & 3.540 & 0.340\\

P2 & 0.001 & 0 & 8 & 0.667 & 0.566 & 3.540 & 0.340\\

P2 & 0.0025 & 0 & 8 & 0.667 & 0.566 & 3.540 & 0.340\\

P2 & 0.005 & 0 & 6 & 0.583 & 0.614 & 4.694 & 0.349\\

P2 & 0.01 & 0 & 4 & 0.500 & 0.646 & 5.616 & 0.358\\

P3 & 0 & 0 & 7 & 0.083 & 0.693 & 19.309 & 0.651\\

P3 & 0.001 & 0 & 4 & 0.083 & 0.672 & 17.170 & 0.653\\

P3 & 0.0025 & 0 & 0 & 0.250 & 0.630 & 12.473 & 0.651\\

P3 & 0.005 & 0 & 0 & 0.250 & 0.630 & 12.473 & 0.651\\

P3 & 0.01 & 0 & 0 & 0.250 & 0.630 & 12.473 & 0.651\\

P4 & 0 & 0 & 7 & 0.500 & 0.577 & 3.970 & 0.467\\

P4 & 0.001 & 0 & 7 & 0.500 & 0.577 & 3.970 & 0.467\\

P4 & 0.0025 & 0 & 4 & 0.667 & 0.559 & 2.814 & 0.460\\

P4 & 0.005 & 0 & 1 & 0.917 & 0.528 & 0.953 & 0.436\\

P4 & 0.01 & 0 & 0 & 1.000 & 0.520 & 0.494 & 0.434\\

P5 & 0 & 0 & 9 & 0.750 & 0.575 & 8.443 & 0.378\\

P5 & 0.001 & 0 & 8 & 0.667 & 0.625 & 11.296 & 0.399\\

P5 & 0.0025 & 0 & 6 & 0.500 & 0.727 & 17.412 & 0.427\\

P5 & 0.005 & 0 & 6 & 0.500 & 0.727 & 17.412 & 0.427\\

P5 & 0.01 & 0 & 3 & 0.250 & 0.879 & 26.373 & 0.474\\

P6 & 0 & 0 & 3 & 0.250 & 0.517 & 12.708 & 0.488\\

P6 & 0.001 & 0 & 1 & 0.083 & 0.556 & 15.705 & 0.521\\

P6 & 0.0025 & 0 & 1 & 0.083 & 0.556 & 15.705 & 0.521\\

P6 & 0.005 & 0 & 0 & 0.000 & 0.575 & 17.128 & 0.528\\

P6 & 0.01 & 0 & 0 & 0.000 & 0.575 & 17.128 & 0.528\\

\bottomrule\end{longtable}\endgroup

\begingroup\small\setlength{\tabcolsep}{3pt}\begin{longtable}{llrrrrrr}\caption{Selector sweep: duration.}\\\toprule

Process & Threshold & D & R & Agree & Cost ratio & Regret & Event\\\midrule\endfirsthead

\toprule Process & Threshold & D & R & Agree & Cost ratio & Regret & Event\\\midrule\endhead

P1 & 0 & 3 & 3 & 0.250 & 0.542 & 2.765 & 0.378\\

P1 & 0.0025 & 1 & 3 & 0.250 & 0.548 & 2.895 & 0.384\\

P1 & 0.005 & 0 & 3 & 0.250 & 0.549 & 2.920 & 0.384\\

P1 & 0.01 & 0 & 3 & 0.250 & 0.549 & 2.920 & 0.384\\

P1 & 0.02 & 0 & 3 & 0.250 & 0.549 & 2.920 & 0.384\\

P2 & 0 & 4 & 6 & 0.583 & 0.596 & 4.266 & 0.344\\

P2 & 0.0025 & 0 & 6 & 0.583 & 0.614 & 4.694 & 0.349\\

P2 & 0.005 & 0 & 6 & 0.583 & 0.614 & 4.694 & 0.349\\

P2 & 0.01 & 0 & 6 & 0.583 & 0.614 & 4.694 & 0.349\\

P2 & 0.02 & 0 & 6 & 0.583 & 0.614 & 4.694 & 0.349\\

P3 & 0 & 5 & 0 & 0.333 & 0.612 & 10.592 & 0.641\\

P3 & 0.0025 & 0 & 0 & 0.250 & 0.630 & 12.473 & 0.651\\

P3 & 0.005 & 0 & 0 & 0.250 & 0.630 & 12.473 & 0.651\\

P3 & 0.01 & 0 & 0 & 0.250 & 0.630 & 12.473 & 0.651\\

P3 & 0.02 & 0 & 0 & 0.250 & 0.630 & 12.473 & 0.651\\

P4 & 0 & 4 & 1 & 0.917 & 0.528 & 0.953 & 0.436\\

P4 & 0.0025 & 1 & 1 & 0.917 & 0.528 & 0.953 & 0.436\\

P4 & 0.005 & 0 & 1 & 0.917 & 0.528 & 0.953 & 0.436\\

P4 & 0.01 & 0 & 1 & 0.917 & 0.528 & 0.953 & 0.436\\

P4 & 0.02 & 0 & 1 & 0.917 & 0.528 & 0.953 & 0.436\\

P5 & 0 & 4 & 6 & 0.500 & 0.726 & 17.393 & 0.427\\

P5 & 0.0025 & 0 & 6 & 0.500 & 0.727 & 17.412 & 0.427\\

P5 & 0.005 & 0 & 6 & 0.500 & 0.727 & 17.412 & 0.427\\

P5 & 0.01 & 0 & 6 & 0.500 & 0.727 & 17.412 & 0.427\\

P5 & 0.02 & 0 & 6 & 0.500 & 0.727 & 17.412 & 0.427\\

P6 & 0 & 2 & 0 & 0.000 & 0.576 & 17.168 & 0.526\\

P6 & 0.0025 & 0 & 0 & 0.000 & 0.575 & 17.128 & 0.528\\

P6 & 0.005 & 0 & 0 & 0.000 & 0.575 & 17.128 & 0.528\\

P6 & 0.01 & 0 & 0 & 0.000 & 0.575 & 17.128 & 0.528\\

P6 & 0.02 & 0 & 0 & 0.000 & 0.575 & 17.128 & 0.528\\

\bottomrule\end{longtable}\endgroup

Raw trajectory archives retain all independent arrivals and candidate action outcomes. Per-fit records include actual thresholds, their relation to calibration support, activation, lead time, false and unnecessary activation, WIP area, completed-case durations and residual work. Completed-case duration is censored by the evaluation horizon; complete-cohort replays, where available, are separate extended-horizon diagnostics. P4 records preserve outcomes for both competing processes.

\clearpage

\subsection{Trigger support and recovery}

The following routing-feedback table reports retained-load choices. Threshold counts show actual WIP values and numbers of fits. The training range is the range, across fits, of maximum observed training WIP; $q/\max w$ is calculated within each fit. K/D denote known/discovered structure. These diagnostics do not alter activation rules.

\begin{center}\small\setlength{\tabcolsep}{3pt}\begin{tabular}{llp{.24\textwidth}rrr}\toprule

Process & N & Threshold:count & Train max & $q/\max w$ & Above\\\midrule

P1 & K & 6:11; 12:1 & 5--8 & 0.75--2.40 & 4/12\\

P1 & D & 6:12 & 5--8 & 0.75--1.20 & 4/12\\

P2 & K & 8:12 & 6--12 & 0.67--1.33 & 2/12\\

P2 & D & 8:12 & 6--12 & 0.67--1.33 & 2/12\\

P3 & K & 10:11; none:1 & 7--11 & 0.91--1.43 & 5/12\\

P3 & D & 10:11; none:1 & 7--11 & 0.91--1.43 & 5/12\\

P4 & K & 10:12 & 11--19 & 0.53--0.91 & 0/12\\

P4 & D & 10:12 & 11--19 & 0.53--0.91 & 0/12\\

P5 & K & 10:11; 20:1 & 12--18 & 0.56--1.18 & 1/12\\

P5 & D & 10:11; 20:1 & 12--18 & 0.56--1.18 & 1/12\\

P6 & K & 10:11; none:1 & 9--13 & 0.77--1.11 & 1/12\\

P6 & D & 10:12 & 9--13 & 0.77--1.11 & 1/12\\

\bottomrule\end{tabular}\end{center}

\begin{center}\small\setlength{\tabcolsep}{4pt}\begin{tabular}{llrrrrr}\toprule

Process & N & Activate & Nonsevere & Lead/$T$ & Event & Prevented\\\midrule

P1 & K & 0.717 & 0.371 & 5.661 & 0.365 & 0.326\\

P1 & D & 0.724 & 0.385 & 5.979 & 0.378 & 0.299\\

P2 & K & 0.736 & 0.542 & 5.666 & 0.325 & 0.223\\

P2 & D & 0.736 & 0.542 & 5.666 & 0.370 & 0.117\\

P3 & K & 0.628 & 0.052 & 4.649 & 0.651 & 0.023\\

P3 & D & 0.628 & 0.052 & 4.649 & 0.651 & 0.023\\

P4 & K & 0.828 & 0.533 & 8.362 & 0.490 & 0.223\\

P4 & D & 0.828 & 0.533 & 8.362 & 0.490 & 0.223\\

P5 & K & 0.833 & 0.640 & 9.790 & 0.354 & 0.330\\

P5 & D & 0.833 & 0.640 & 9.790 & 0.354 & 0.330\\

P6 & K & 0.667 & 0.203 & 6.346 & 0.401 & 0.382\\

P6 & D & 0.729 & 0.227 & 6.302 & 0.528 & 0.187\\

\bottomrule\end{tabular}\end{center}

Nonsevere activation is conditional on no severe event in matched no action; it is not necessarily a harmful intervention. Lead time is measured against that counterfactual event onset. Prevention is conditional on an event without action. Recovery below is WIP at most $2S$ for $10T$ after a pulse, censored at the original $180T$ horizon. Entries average within each independent fit before averaging over fits; they concern all pulse trajectories, not only crisis trajectories.

\begin{center}\small\begin{tabular}{lrrrr}\toprule

Process & No action & Predictive K & Retained K & Generator\\\midrule

P1 & 0.792 & 1.000 & 1.000 & 1.000\\

P2 & 0.742 & 0.958 & 1.000 & 1.000\\

P3 & 0.169 & 0.667 & 0.560 & 0.727\\

P4 & 0.555 & 0.979 & 0.878 & 0.992\\

P5 & 0.815 & 0.917 & 0.987 & 1.000\\

P6 & 0.477 & 0.935 & 0.953 & 1.000\\

\bottomrule\end{tabular}\end{center}

\subsection{Shared-resource consequences and bottleneck location}

The resource diagnostic replays P4/P5 under all four mechanisms and all 12 fits, using the two high-baseline pulse conditions and replicate 0. Every action type at the smallest frozen threshold is included, together with saved selected policies. All 880 prefixes through $180T$ exactly match original test samples. Arrivals then stop and actions remain active through $600T$ to complete the arrival cohort. Queue integrals use only the original horizon. Complete-cohort duration therefore has a different horizon from the primary endpoint.

\begin{longtable}{llrrrr}\caption{P4: process-specific WIP integrals (thousands) from the full test, and complete-cohort cycle times from the extended diagnostic.}\\\toprule

Mechanism & Action at 10 & WIP X & WIP Y & Cycle X & Cycle Y\\\midrule\endfirsthead

\toprule Mechanism & Action at 10 & WIP X & WIP Y & Cycle X & Cycle Y\\\midrule\endhead

fixed & None & 7.78 & 8.61 & 21.53 & 23.86\\

fixed & Capacity & 4.68 & 5.41 & 11.94 & 13.80\\

fixed & Rework X & 7.44 & 8.45 & 20.61 & 23.36\\

fixed & Rework Y & 7.46 & 7.99 & 20.59 & 22.05\\

service & None & 43.69 & 45.54 & 184.34 & 191.92\\

service & Capacity & 22.81 & 24.08 & 84.37 & 87.91\\

service & Rework X & 41.85 & 44.62 & 174.80 & 186.61\\

service & Rework Y & 42.17 & 42.49 & 174.97 & 174.70\\

routing & None & 31.13 & 30.37 & 111.82 & 109.49\\

routing & Capacity & 15.01 & 14.88 & 46.73 & 46.50\\

routing & Rework X & 16.54 & 21.86 & 54.25 & 72.12\\

routing & Rework Y & 21.54 & 16.27 & 70.68 & 52.47\\

combined & None & 91.64 & 89.79 & 567.45 & 554.47\\

combined & Capacity & 72.57 & 70.68 & 373.43 & 365.29\\

combined & Rework X & 69.66 & 83.69 & 355.70 & 471.85\\

combined & Rework Y & 86.18 & 69.81 & 491.10 & 366.32\\

\bottomrule\end{longtable}

\begin{longtable}{llrrr}\caption{P5: mean queue lengths at B and D during the original horizon; number of paired paths whose largest queue moves from B to D.}\\\toprule

Mechanism & Action at 10 & Queue B & Queue D & B to D\\\midrule\endfirsthead

\toprule Mechanism & Action at 10 & Queue B & Queue D & B to D\\\midrule\endhead

fixed & None & 13.12 & 0.11 & 0/24\\

fixed & Capacity B & 1.99 & 4.75 & 24/24\\

fixed & Capacity D & 13.12 & 0.01 & 0/24\\

fixed & Rework & 13.12 & 0.08 & 0/24\\

service & None & 140.94 & 0.01 & 0/24\\

service & Capacity B & 38.36 & 0.09 & 0/24\\

service & Capacity D & 140.94 & 0.00 & 0/24\\

service & Rework & 140.94 & 0.00 & 0/24\\

routing & None & 13.28 & 35.73 & 0/24\\

routing & Capacity B & 2.09 & 46.05 & 0/24\\

routing & Capacity D & 13.28 & 4.24 & 0/24\\

routing & Rework & 13.28 & 1.69 & 0/24\\

combined & None & 141.02 & 1.49 & 0/24\\

combined & Capacity B & 40.19 & 118.56 & 24/24\\

combined & Capacity D & 141.02 & 0.05 & 0/24\\

combined & Rework & 141.02 & 0.03 & 0/24\\

\bottomrule\end{longtable}


\label{maintextend}
\bibliographystyle{splncs04}
\bibliography{references}
\end{document}
